%%%%%%%%%%%%%%%%%%%%%%%%%%%%%%%%%
%
%       EXERCÍCIO
%
%%%%%%%%%%%%%%%%%%%%%%%%%%%%%%%%%

\ifdefstring{\atividade}{prova}{%
    \ifdefstring{\modo}{objetivo}{%
        \renewcommand{\valorquestao}{\ValorQObj\ ponto}
    }{%
        \renewcommand{\valorquestao}{\ValorQDisc\ pontos}
    }%
}%

\begin{exercicioBanco}[\valorquestao]
No universo \(\bbN = \{1,2,3,\ldots\}\), sejam:

\begin{enumerate}[\(\bullet\)]
    \begin{multicols}{2}
        \item \(A\) o conjunto dos números múltiplos de \(4\);
        \item \(B\) o conjunto dos números múltiplos de \(3\);
        \item \(C\) o conjunto dos números múltiplos de \(6\).
    \end{multicols}
\end{enumerate}

Determine os \(4\) menores números que pertencem ao conjunto

\[
    B\setminus(A \cup C).
\]

% Define as alternativas
\newcommand{\alternativas}{%
    \begin{center}
        \begin{tabularx}{\textwidth}{XXXXX}
            (a) \(3, 6, 9, 12\). &
            (b) \(3, 9, 15, 21\). &
            (c) \(3, 9, 18, 21\). &
            (d) \(6, 9, 12, 15\). &
            (e) \(3, 6, 15, 18\).
        \end{tabularx}
    \end{center}
}

% Define a resposta correta
\newcommand{\resposta}{B}

% Lógica condicional para exibição
\ifdefstring{\atividade}{lista}{%
    \alternativas
    \vspace{0.5em}
    
    \noindent\textbf{Resposta:} letra \textbf{\resposta}.
}{%
    \ifdefstring{\modo}{objetivo}{%
        \alternativas
    }{%
        % modo = discursiva → não mostra alternativas
    }
}
\end{exercicioBanco}

